
\section{Related Work}

Terrain generation methods can be broadly classified into three families: procedural and simulation-based models, sketch-based control techniques, and learning-based synthesis. Our approach draws from all three, combining procedural control with learning-based realism.

\paragraph{Noise- and simulation-based terrains}
Procedural noise functions such as Perlin and simplex noise, and fractal Brownian motion (fBm), have long been used to synthesize diverse terrains efficiently~\cite{Perlin1985,Perlin2001,Musgrave1989,Ebert2003,Reinhard2010}. Although these methods provide scalability and interactivity, they lack explicit correspondence to geological or biological processes, limiting their realism for coral reef morphologies.  
Physical and simulation-based models incorporate erosion, uplift, or sediment transport to achieve higher realism~\cite{Benes2006,Mei2007,Neidhold2005,Cordonnier2016,Cordonnier2017a,Schott2023}. However, these simulations are computationally expensive and are typically calibrated for continental or fluvial terrains rather than shallow marine systems. As a result, neither class of approaches directly captures the coupled volcanic, subsidence, and coral growth phenomena characteristic of reef islands.

\paragraph{Sketch-based control}
To increase artistic control and interactivity, sketch-based methods allow users to specify structural constraints or silhouettes directly~\cite{Gain2009,Hnaidi2010,Tasse2014,Gasch2020,Talgorn2018,Guerin2022}. Curve and gradient-domain formulations enable intuitive authoring of shapes and elevation fields while maintaining spatial coherence. These approaches are particularly relevant for procedural workflows because they reduce parameter tuning and expose higher-level control to non-expert users.  
Our method extends this paradigm through a dual-view sketch interface, combining top and profile sketches, with stroke-defined deformation fields while preserving semantic labels that later condition the learning stage.

\paragraph{Learning-based terrains}
Data-driven methods learn terrain appearance from examples, often through adversarial or diffusion-based generative models. Generative adversarial networks (GANs) and their conditional variants (cGANs) have been used to synthesize height fields and realistic landscapes~\cite{Goodfellow2014,WulffJensen2018,Spick2019,Guerin2017,Panagiotou2020,Beckham2017}. Among these, pix2pix~\cite{Isola2017} remains an effective baseline for translating semantic maps to terrain elevations. Recent diffusion models~\cite{Ho2020,Nichol2021,Lochner2023} achieve improved detail and realism but at significantly higher computational cost, making them less suitable for real-time interaction.  
In contrast, our approach relies on a lightweight pix2pix architecture trained entirely on synthetic data produced by the procedural sketch model, bridging controllable authoring and learned realism.

% \medskip
In summary, previous methods either offer control with limited realism (procedural/sketch-based) or realism without control (learning-based). Our work unifies both by procedurally encoding geological and biological priors into a dataset for a conditional generative model capable of real-time, user-driven synthesis of coral reef islands.






% \section{State of the art}
% \label{sec:coral-island-SotA}

% Procedural terrain generation and sketch-based modelling have each seen significant advances over the past decades, yet neither alone fully addresses the particular challenges of coral reef island synthesis. On the one hand, classic noise-based and physically driven methods (Perlin noise, hydraulic erosion, uplift-erosion models) excel at producing broad, natural-looking landscapes but lack the biologically inspired reef geometries and dynamic island evolution governed by subsidence and coral growth. On the other hand, sketch-based tools give users intuitive control over silhouettes and terrain profiles, but typically require expert parameter tuning to achieve realism and do not model long-term geological or ecological processes. More recently, deep learning, especially GANs and cGANs, has emerged as a powerful way to learn terrain features and geological rules from data, yet it relies on large, labelled datasets that are scarce for coral reef islands. 

% In this section, we first review the key geological theories that explain reef formation, then examine traditional procedural terrain generators, followed by an overview of sketch-based terrain editing frameworks, and finally discuss recent advances in GAN-based terrain synthesis.

% \subsection{Coral reef island formation theories}
% \label{sec:coral-island-sota-coral-theories}

% Since Darwin first proposed his subsidence theory, presented in \cref{sec:coral-island-introduction}, geologists and biologists have debated how exactly coral reefs have formed around land masses. In this section we present three major alternatives to Darwin's theory, before explaining why Darwin's unified subsidence model provides the most direct foundation for our procedural generation.

% \begin{Itemize}
%     \AltTextImageBottom{
%         \Item{} \textbf{Murray's stand-still theory \cite{Murray1880}} \\
%         Murray argued that reefs could develop on stable, non-sinking platforms, with coral growth keeping pace with modest sea-level changes rather than underlying land subsidence. He proposed that as long as water depth remained within the photic zone, reefs would accrete upward solely in response to environmental sea-level fluctuations, and continuously seaward, forming a gradual structure from a fringing reef to barrier reef (\cref{fig:coral-island-murray-theory}). Atolls in this theory are mainly created by the degradation of the middle island through aeolian and coastal erosion until it becomes completely submeged. Murray's emphasis on environmental stability and gradual sea-level change was supported by early observations of island terraces and reef growth patterns. However, later drilling and seismic data revealed volcanic foundations buried beneath reef limestones on many atolls, which are evidence of true subsidence that Murray's theory cannot explain.
%     }{Murray.pdf}% }{Murray-drawing1.png, Murray-drawing2.png, Murray-drawing3.png}
%     {In Murray's theory, with the island fixed in place but reducing in size due to erosion, reefs (red) simply remains at sea level and grows outwards (producing fringing reefs, then barrier reefs, and eventually atolls) without any subsidence, leaving only the inside reef filled with sand (yellow) forming a lagoon.}{fig:coral-island-murray-theory}{Murray's stand-still theory}

%     \AltTextImageBottom{
%         \Item{} \textbf{Daly's glacial-control theory \cite{Daly1915}} \\  
%         Daly shifts the focus from steady subsidence to global sea-level oscillations driven by glacial-interglacial cycles. He argues that reefs thrive during high sea levels in interglacials but are exposed and eroded during glacial lowstands. As sea levels rise again, the wave-cut benches left behind provide ideal habitats for new coral colonies, which grow upward with the water level and can develop into barrier reefs or atolls once sea level overtops the island (\cref{fig:coral-island-daly-theory}). He supports this model by pointing to multiple reef terraces at different elevations and well-documented 100m sea-level drops during ice ages as support for this cycle-driven reef development. The plausibility of Daly's model lies in its direct link with climatic data and the clear geomorphic signatures of past sea-level stands. Yet, core samples often show uninterrupted reef accretion atop subsiding volcanic bases, indicating that glacial cycles alone cannot account for continuous reef "keep-up" growth.
%     }{Daly.pdf} % }{Daly-drawing1.png, Daly-drawing2.png, Daly-drawing3.png}
%     {In Daly's theory, repeated sea-level drops expose the reef to wave planation and erosion into a flat bench (grey), and subsequent high stands see new coral rims (red) accrete on that outer terrace to form a lagoon-separated barrier reef.}{fig:coral-island-daly-theory}{Daly's glacial-control theory}

%     \AltTextImageBottom{
%         \Item{} \textbf{Droxler's karstification theory \cite{Droxler2021}} \\  
%         In contrast to models based on volcanic subsidence, Droxler et Jorry attribute atoll formation to repeated karst dissolution of a broad carbonate platform during low-sea-level glacials, followed by renewed coral accretion in interglacials (\cref{fig:coral-island-droxler-theory}). High-resolution bathymetry and drill cores reveal karst-etched depressions beneath certain atoll crests, supporting this cyclic exposure-dissolution mechanism. While compelling for extensive carbonate shelves, this theory hinges on pre-existing flat platforms and meteoric water circulation (water originating from precipitation), which are conditions uncommon on volcanic seamounts. As a result, Droxler's model explains atoll rings on continental carbonate foundations but does not readily apply to volcanic island reef systems. 
%     }{Droxler.pdf} % }{Droxler_placeholder.png}
%     {A broad carbonate platform (grey) is exposed and karstified during low-sea-level glacials, then drowns and is rimmed by reef growth (red) at its outer edge during high stands, yielding an atoll-style rings.}{fig:coral-island-droxler-theory}{Droxler's karstification theory}

% \end{Itemize}

% Of all the competing hypotheses, Charles Darwin's subsidence model offers several compelling advantages for our procedural island generation due to its simplicity, historical significance, and effectiveness \cite{Tomascik1997}. This theory provides a straightforward framework outlining a clear progression from fringing reefs to barrier reefs to atolls as a volcanic island subsides, which can be easily modelled, making it practical for generating plausible island landscapes. This simplicity not only ensures predictability in simulation outcomes but also reduces computational demands, which is particularly beneficial for interactive edition applications. Additionally, Darwin's model provides a foundational basis upon which more complex phenomena (sea-level changes and climatic effects), considered as secondary factors in the geological formation of islands, could be layered. This allows us to start with a basic model and, in the future, add complexity as needed, offering flexibility in simulation design.

% In our approach, we translate the core principles of Darwin's theory into a procedural generation model, emulating the gradual sinking of volcanic islands while coral reefs grow to keep pace with changing sea levels. This allows us to realistically model the progressive transformation of islands from volcanic landmasses to coral-dominated atolls. By capturing this interplay, we procedurally generate a wide variety of island structures that reflect real-world geological processes.

% In generating synthetic coral reef islands, we adopt a set of simplifying assumptions that are interpreted by the formation process from Darwin's theory and field observations:
% \begin{Itemize}
%     \AltTextImage{
%         \Item{} \textbf{Radial symmetry and radial arrangement of features} \\
%         While volcanic islands do not have a perfect circular shape and reefs are not exact rings, they grow outward in roughly concentric belts (looking at the island in \cref{fig:coral-island-radial-symmetry-assumption}, we see the typical layout mountain-beach-lagoon-reef-ocean). Anchoring our sketch-to-terrain engine on this radially organised structure offers two main benefits to the user: a single intuitive control over the overall island shape and while keeping distance computations simple and efficient.

%         \Item{} \textbf{Uniform profile shape} \\
%         Although real islands have ridges, valleys, local bumps and many other features, we simplify the model to a single elevation curve from centre to edge, making the profile-view sketch a one-dimensional drawing, enabling users to modify elevation with a single stroke.
%     }
%     {Tuvuca_island.png}{Aerial view of Tuvuca Island, Lau Archipelago, Fiji.}{fig:coral-island-radial-symmetry-assumption}{Aerial view of Tuvuca Island}

%     \AltTextImage{
%         \Item{} \textbf{Subsidence and coral "keep up" growth} \\
%         Actual reefs respond to light, nutrients, and biology, but the dominant effect is that corals grow vertically to stay near the photic zone. We decouple coral growth from subsidence and simply keep reef heights in a fixed band beneath the surface, capturing Darwin's core insight using minimal parameters.

%         \Item{} \textbf{Wind and wave deformation} \\
%         To represent coastal erosion resulting from sediment transport and complex hydrodynamics \cite{Terry2013}, we simplify the model by using a vector field painted by the user representing wind and wave directions, deforming the global island (\cref{fig:coral-island-wind-deformation-assumption} suggest that island outlines are influenced by the average wind direction). This vector field allows radial symmetry to be broken in a controllable way, without relying on a full erosion simulator.
%     }
%     {Mauritius_wind.png}{Mauritius is deformed. We see a deep slope on the west side while the east side has a gentle slope. The wind rose, showing in which direction the mean wind velocity is pointing, is correlated with this deformation.}{fig:coral-island-wind-deformation-assumption}{Mauritus Island and average wind direction}

%     \Item{} \textbf{Independence of islands} \\
%     While archipelagos share currents, sediments, and flood each other's lagoons, simulating multi-island interactions is computationally expensive and hard to control. We treat each island as an isolated entity to generate multi-island scenes by using a simple blending of individual islands, keeping our method fast and predictable.

% \end{Itemize}

% \midConclusion

% These assumptions, grounded in theories of coral reef island formation, provide a practical foundation for procedural modelling. While they simplify the full complexity of geological and biological processes, they capture the essential dynamics needed to produce plausible island structures. Simultaneously, they ensure our system remains intuitive, controllable, and computationally efficient, enabling the generation of diverse and plausible islands suitable for interactive design.





% \subsection{Traditional terrain generation methods}
% \label{sec:coral-island-sota-traditional}

% Procedural generation of terrain has been a well-researched area in computer graphics and simulations, where the goal is to create large, realistic landscapes with minimal manual input. Various methods have been developed over the years to generate terrains automatically, from noise-based methods to physically based erosion simulations, sketch-driven mechanisms, and more recently, deep learning techniques.

% However, coral reef islands present unique challenges due to the combination of long-term geological processes (mainly subsidence and coral reef growth) and environmental interactions (i.e. erosion caused by wind and waves). In this section, we review the key techniques proposed for terrain generation, highlighting their limitations for coral reef island formation, and positioning our work as an approach that addresses these challenges.

% \subsubsection{Noise-based terrain generation}

% \begin{figure}
%     \includegraphics[width = \linewidth]{noise_examples3.pdf}
%     \caption[Islands generation using fBm noise and falloff masks]{Three different results of an island generated from noise functions. In each case, the initial height field is identical, computed using flow noise. The falloff masks are created by combining fBm noise, modulated by the Euclidean distance from the image centre, and further shaped through domain warping and power-law remapping of the mask values. The only parameter modified between the three examples is the exponent used in this remapping. The resulting islands differ significantly in shape and scale, and the method offers limited control over specific coastal features such as lagoons and reefs.}
%     \label{fig:coral-island-noise-example}
% \end{figure}

% Noise-based procedural generation remains one of the most widely used techniques for creating natural-looking terrains on height fields. Perlin noise \cite{Perlin1985}, Simplex noise \cite{Perlin2001}, and the Diamond-square algorithm \cite{Fournier1982} are foundational algorithms that generate pseudo-random yet continuous variations across a grid, producing terrain features that resemble organic landscapes. These techniques have been widely adopted in computer graphics and game development due to their efficiency and visual appeal.

% Beyond basic noise functions, advanced techniques such as fractal Brownian motion (fBm) and multifractal noise have been introduced to add finer-scale variation and detail \cite{Musgrave1989,Ebert2003}. fBm combines multiple layers, or "octaves," of noise at different frequencies and amplitudes, producing terrains that exhibit realistic and varied features. The combination of noise with domain warping and signal processing techniques has been explored in depth in procedural modelling literature \cite{Reinhard2010}, enabling further control over visual complexity and terrain realism.

% Noise functions are often paired with falloff maps to produce island-like terrains, where elevation gradually decreases towards the edges of the domain, mimicking coastlines and basic island shapes (see \cref{fig:coral-island-noise-example}). Various methods for enhancing island generation using noise and falloff blending have been proposed for applications in games and virtual worlds \cite{Olsen2004}. While these techniques excel at producing large, visually diverse landscapes quickly, they suffer from several key limitations when applied to the modelling of coral reef islands.

% However, limitations arise when applying these methods to coral reef islands as noise-based modelling lack grounding in geological or biological reality. They generate spatial patterns through mathematical noise, not through simulations of real-world processes such as volcanic subsidence or coral accretion. The signal processing parameters typically involved (frequency, lacunarity, gain, amplitude, ...) are tuned for visual effect rather than scientific plausibility. As highlighted in procedural modelling surveys \cite{Smelik2009,Galin2019}, this disconnect results in a lack of semantic control and poor correlation with actual environmental dynamics, making it difficult to represent phenomena such as reef rings, lagoons, or atoll structures in a biologically or geologically coherent way.

% Moreover, the biological aspects of coral growth are inherently tied to environmental conditions. Coral reefs form and persist only within specific ranges of water depth, sunlight, salinity, and water quality. Their growth patterns are further influenced by ecological health, nutrient availability, and symbiotic relationships. These dependencies are extremely difficult to capture in procedural noise models, which are not designed to reproduce such complex and coupled dynamics.

% Our approach goes beyond the randomness of noise-based generation by incorporating real-world geological and biological processes into the terrain formation pipeline. Specifically, we model the gradual subsidence of volcanic islands and the upward growth of coral reefs, both of which are central to the long-term evolution of coral reef islands. By embedding these natural processes directly into the generation algorithm, we produce terrains that are not only more realistic but also more controllable. This integration of scientific modelling with procedural flexibility allows us to overcome the inherent limitations of traditional noise-based techniques and more accurately represent the complex formation of coral reef island systems.

% \subsubsection{Simulation-based modelling}

% Simulation-based terrain modelling methods aim to increase realism by replicating natural processes such as erosion, sediment transport, tectonic uplift, and vegetation growth. Unlike noise-based strategies, which rely on random functions, simulation-based processes model causality and temporal dynamics to describe how a terrain evolves over time under physical or biological forces. These methods are often used to enhance base terrains, adding geologically plausible detail and structure \cite{Benes2006, Smelik2009}.

% \subsubsubsection{Hydraulic and thermal erosion}

% Hydraulic erosion models simulate the impact of flowing water on the landscape by modelling erosion, sediment pickup, transport, and deposition. Early implementations by \cite{Musgrave1989} laid the groundwork for erosion in procedural generation, while more recent works have accelerated these simulations using GPU architectures \cite{Mei2007} and particle-based methods \cite{Neidhold2005}. These simulations either follow Eulerian fluid models or Lagrangian particle systems to capture terrain displacement.

% Thermal erosion, by contrast, simulates mass movement due to gravity, redistributing material from steeper slopes to gentler gradients, akin to landslides or soil creep \cite{Benes2006}. These erosion models generate realistic fluvial networks and landforms, but they are parameter-sensitive and computationally expensive.

% Moreover, such models are generally designed for terrestrial landscapes and lack mechanisms for simulating underwater sedimentation, reef growth, or biogenic processes crucial to coral reef island formation. These models typically simulate time scales relevant to geomorphological processes (hundreds to thousands of years), which are mismatched with both the faster dynamics of biological processes (e.g., coral growth) and the slower geological evolution of reef islands.

% \subsubsubsection{Tectonic uplift and geologic simulation}

% Geological simulation approaches such as those proposed by \cite{Cordonnier2016, Cordonnier2017a} and extended by \cite{Schott2023} model terrain evolution through crustal deformation and tectonic uplift. These methods simulate isostatic adjustments, plate tectonics, or local uplift phenomena, often over geological timescales.

% Although well-suited for mountain-building processes or fault line modelling, these methods are not designed to account for biogenic terrain formation, such as coral reef accretion, which is critical for simulating coral reef islands. As a result, despite being physically grounded models, they do not capture the coupled geological and biological dynamics necessary for representing the long-term evolution of reef islands.

% \subsubsubsection{Vegetation and ecosystem dynamics}

% A number of simulation-based terrain models integrate ecological dynamics to reflect the feedback between terrain and living systems. For instance, \cite{Ecormier-Nocca2021} and \cite{Cordonnier2017b} simulate interactions between vegetation and terrain erosion, modelling plant colonisation, growth, and their influence on soil stability and moisture retention.

% These ecosystem simulations allow more complex landscape evolution by considering biotic agents; however, they are designed primarily for terrestrial plants and temperate ecosystems. Coral colonies, in contrast, are marine organisms with strict environmental requirements (e.g., limited depth, adequate sunlight, nutrients presence, warm water temperatures, ...). Accurately simulating these dependencies would require significant computational resources. 

% Furthermore, coral growth is not a passive process such as sediment accumulation, but an active accretion system that builds calcium carbonate structures over thousands of years. These unique growth mechanisms, constrained by marine ecology, fall outside the scope of existing vegetation or soil-plant-water feedback models.

% \citeName{Li2021} propose to procedurally synthesizes coral colonies structures by emulating stochastic reef growth with a Diffusion-Limited Aggregation process, producing reef patterns but does not propose control over the simulation \cite{Li2021}. The integration of user input such as sketch interfaces remains essential for controlling the generated output.

% \midConclusion

% While simulation-based models represent a significant advancement over purely procedural models, they fall short in capturing the coupled geological and biological dynamics that shape coral reef islands. They are either computationally intensive, domain-specific, or biologically inapplicable, highlighting the need for a new class of terrain generation tools that embed long-term marine biogeomorphological processes into the procedural pipeline.


% \subsection{Sketch-based terrain modelling}
% \label{sec:coral-island-sota-sketches}

% The term "sketching" encompasses several definitions: either refering to performing hand or body gestures, creating a rough drawing, or outlining an idea in a simplified form. Accordingly, sketch-based modelling in 3D computer graphics can be understood through three complementary perspectives, each centred around a distinct core concept: interaction, construction and interpretation.

% First, "sketching" may focus on interaction, where gestures captured through hand or body motion are used to manipulate virtual objects, often in immersive environments (virtual or augmented reality), using established techniques such as sculpting and distortion \cite{Olsen2009, Cook2009}. Second, "sketching" may involve construction, where simple geometric primitives (curves, parametric shapes, implicit surfaces, ...) are combined under constraints to build a more complex model. Finally, "sketching" may centre on interpretation, where the user draws strokes on a 2D canvas and the system analyses their meaning to generate a plausible 3D model.

% While sketch-based modelling encompasses a wide range of techniques, including gesture-driven interaction in immersive environments, this work focuses primarily on construction and interpretation aspects. In particular, construction serves as the foundation for procedural generation techniques using geometric primitives and constraints, while interpretation becomes relevant when exploring data-driven approaches using deep learning to infer terrain structure from sketches. Interaction-based techniques, though significant in other contexts, fall outside the scope of this work.

% To clarify terminology in this chapter, we refer to the constructive approach as sketch-based, and to the interpretive, learning-driven approach as learning-based. It is important to note, however, that the boundaries between these categories are inherently blurry and often overlap in practice.

% In procedural terrain generation, sketch-based construction workflows enable users to shape landscapes by manipulating high-level geometric primitives through intuitive sketching interfaces. These methods allow the definition of key terrain features (mountains, valleys, coastlines, ...) by drawing their outlines on a two-dimensional canvas, which are then procedurally transformed into 2.5D or 3D terrain representations. This technique offers a high degree of control, making it particularly effective for creative applications such as video games and robotics simulations, where modelling is primarily user-driven.

% \subsubsection{Curve-based modelling}

% Sketch-based terrain generation often begins with user-defined curves which act as high-level constraints to guide the shape of the terrain. These curves represent silhouettes, ridgelines, valleys, or feature outlines. Once defined, they are interpreted by the system and translated into elevation changes through various computational techniques. This type of approach allows for intuitive control over large-scale landforms while maintaining a procedural foundation for terrain synthesis.

% \AltTextImage{
%     One of the earliest and most influential works in this domain was introduced by \citeName{Gain2009} (\cref{fig:coral-island-Gain2009}), enabling users to sketch silhouettes, ridges, and spine curves to define complex terrain structures \cite{Gain2009}. The method employs multiresolution surface deformation and propagates wavelet-based noise from the sketched features to their surroundings, allowing users to generate detailed, natural-looking terrains from minimal input. This method demonstrated the effectiveness of combining intuitive sketch input with procedural detail synthesis.

%     We draw direct inspiration from this work's dual-view sketching strategy, combining top-view and profile sketches, which closely aligns with our concentric curve and height-profile input approach.
% }{sketchingGain2009-s.png}{\citeName{Gain2009} propose to edit terrains through sketching by diffusing a noisy height function over a bounded region (left), effectively modelling mountain ranges (right) \cite{Gain2009}.}{fig:coral-island-Gain2009}{\cite{Gain2009}'s sketching method}

% \AltTextImage{
%     Building on this idea, \citeName{Hnaidi2010} proposed a technique based on diffusion equations \cite{Hnaidi2010}. In their method, curves are annotated with geometric constraints (elevation, slope, and roughness), and a diffusion process is used to interpolate these constraints across the terrain surface. This results in smooth, continuous elevation fields that conform to user-defined features such as rivers, ridgelines, or cliffs (\cref{fig:coral-island-Hnaidi2010}). The use of parameterised curves as terrain anchors allows for precise control over landform shaping, while maintaining a high degree of automation.
% }{sketching-Hnaidi2010.png}{\citeName{Hnaidi2010} define the surface by elevation curves with constrained slopes, forming with ease river beds and cliffs \cite{Hnaidi2010}.}{fig:coral-island-Hnaidi2010}{\cite{Hnaidi2010}'s sketching method}

% \AltTextImage{
%     In a different interaction paradigm, \citeName{Tasse2014} introduced a first-person sketching interface, where users draw terrain silhouettes from a particular camera viewpoint \cite{Tasse2014}. These silhouettes are then projected into 3D, and a deformation algorithm adjusts the terrain so that the drawn features are visible exactly as intended from the user's perspective (\cref{fig:coral-island-Tasse2014}). This method supports complex silhouettes with occlusions, T-junctions, and cusps, and represents an immersive and perceptually grounded approach to sketch-based terrain editing.
% }{sketchingTasse2014-s.png}{\citeName{Tasse2014} use sketching in association with camera position and direction to constrain the edition of height fields \cite{Tasse2014}.}{fig:coral-island-Tasse2014}{\cite{Tasse2014}'s sketching method}

% \smallConclusion

% These methods demonstrate the expressive power of curves as terrain-defining elements. By enabling users to sketch intuitive shapes and constraints, they bridge the gap between artistic intent and procedural complexity. Curve-driven methods remain foundational in terrain modelling, particularly when user control over large-scale structure is essential. While we do not use the diffusion model, the idea of sketch-defined elevation constraints along curves informs our use of user-defined shape boundaries.

% \subsubsection{Constraint-based modelling}

% % While curve-driven techniques provide intuitive shape design, c
% Constraint-based and gradient-based families of methods focus on exerting precise control over terrain features via formal specifications such as elevation values, slopes, or gradient fields. These methods prioritise structural accuracy and procedural consistency, making them particularly suited to applications that demand terrain realism, integration with geographic data, or fine-grained editing capabilities.

% \AltTextImage{
%     A representative example of constraint-based modelling is presented by \citeName{Gasch2020}, proposing a method for procedural terrain generation that respects user-defined elevation constraints \cite{Gasch2020}. Their system allows users to fix values at specific control points (e.g., paths, landmarks) and then solve a system of equations to propagate these constraints throughout the terrain. The method integrates these constraints with a noise-based procedural function to preserve natural randomness while conforming to user intent (\cref{fig:coral-island-Gasch2020}). This approach is especially valuable when generating terrains that must align with real-world data or gameplay constraints.

%     We do not adopt this constraint-solving mechanism, but conceptually relate our profile sketch input to a localised height constraint.

% }{sketching-Gasch2020-s.png}{\citeName{Gasch2020} use a top-view sketch to apply constraints (a) on the noise function modelling a terrain, creating an intuitive means to create paths and roads in the resulting terrain (b and c) \cite{Gasch2020}.}{fig:coral-island-Gasch2020}{\cite{Gasch2020}'s sketching method}

% \AltTextImage{
%     Extending this idea, \citeName{Talgorn2018} introduce a real-time sketch-based terrain generation system based on a GPU-accelerated implementation of midpoint displacement \cite{Talgorn2018}. Users sketch curves with explicit elevation values, which act as absolute constraints, and the system extrapolates and interpolates terrain surfaces in real-time. Moreover, the user input influences the elevation constraint and interpolation method properties of the midpoint displacement algorithm using the Red, Green and Blue channels (\cref{fig:coral-island-Talgorn2018}). Their model supports both global and local control over interpolation curvature and roughness, and introduces semantic labelling of sketched features (e.g., ridges vs. rivers) to influence how terrain propagates around constraints. This combination of sketch-based input, constraint propagation, and semantic control enables expressive, large-scale terrain modelling at interactive speeds. We adopt their notion of semantic labels and hierarchical constraint propagation to support real-time terrain shaping with sketch-defined features without the fractal interpolation.
% }{sketching-Talgorn2018.png}{\citeName{Talgorn2018} use top-view sketching to encode properties in the noise function used to model the terrain surface \cite{Talgorn2018}.}{fig:coral-island-Talgorn2018}{\cite{Talgorn2018}'s sketching method}

% \AltTextImage{
%     Building on the need for more intuitive editing, \citeName{Guerin2022} introduce gradient domain paradigm for terrain modelling \cite{Guerin2022}. Rather than specifying elevation values directly, users interact with slope-based representations, allowing for the manipulation of terrain inclination and the integration of local edits into global terrain structure (\cref{fig:coral-island-Guerin2022}). By controlling terrain gradients and reconstructing elevation through integration, this method enables seamless blending between regions and supports a natural editing workflow, particularly for sculpting realistic mountain ridges, valleys, or plateaus.
%     This gradient-domain editing paradigm offers interesting insights; we adopt an alternative approach as we operate in the elevation domain with semantic control.
% }{sketching-Guerin2022.png}{\citeName{Guerin2022} apply gradient (left) and height (right) constraints on curves and diffuse them on the terrain to create landscapes with lightweight and multi-scale representations \cite{Guerin2022}.}{fig:coral-island-Guerin2022}{\cite{Guerin2022}'s gradient domain terrain modelling}

% \smallConclusion

% Both paradigms offer complementary strengths: constraint-based methods ensure precise adherence to user-defined features or data sources, while gradient-based systems provide fluid, perceptual control over terrain shaping. Together, they represent a shift toward high-level modelling tools that maintain procedural expressiveness while granting users a deeper degree of terrain control.

% \subsubsection{Semantic terrain representation}

% Beyond geometric sketching and low-level constraints, a third class of methods explores high-level terrain construction, where users guide terrain generation using abstract or semantic inputs. These approaches aim to simplify the authoring process by allowing users to describe what a terrain should contain (e.g., a mountain or a valley) without specifying how to generate its geomety. Such methods often rely on symbolic sketching, sparse representations, or domain-specific visual cues, offering powerful tools for conceptual design and inverse procedural modelling.

% \AltTextImage{
%     In this vein, \citeName{Genevaux2015} propose a method for representing terrains as sparse combinations of procedural primitives, referred to as "terrain atoms" \cite{Genevaux2015}. These atoms are stored in a dictionary and are either extracted from real-world data or generated synthetically. The terrain is modelled as a linear combination of these features, forming a Sparse Construction Tree that blends primitives in a compact and expressive form (\cref{fig:coral-island-Genevaux2015}). This representation facilitates terrain editing, amplification, and reconstruction from coarse user input.
% % While this work introduces a symbolic representation of terrain via atoms, it is not used in our method, which instead relies on semantic label maps.
% }{sketching-Genevaux2015-1.png,sketching-Genevaux2015-2.png}{\citeName{Genevaux2015} symbolically define a terrain by an aggregation of patches with real-world meaning \cite{Genevaux2015}.}{fig:coral-island-Genevaux2015}{\cite{Genevaux2015} sparse terrain representation}

% \AltTextImage{
%     An illustrative and domain-specific use case is presented by \citeName{Natali2012}, who introduce a system for rapid visualisation of geological concepts \cite{Natali2012}. Here, users sketch schematic representations of subsurface structures such as faults, folds, or strata, and the system generates plausible 3D visualisations of geological terrains (\cref{fig:coral-island-Natali2012}). Their tool is designed primarily for educational and exploratory purposes, enabling geoscientists and students to create, manipulate, and communicate complex geological scenarios through intuitive sketch input. Although it extends beyond traditional terrain elevation modelling, the work exemplifies how sketch-based systems can operate on a conceptual level and support domain-specific semantics.
%     This work inspired our use of sketch strokes to define deformation fields, although our implementation targets structured terrain generation rather than schematic visualisation.
% }{sketching-Natali-2012-original.png}{\citeName{Natali2012} use sketching to describe geological phenomena: (top) geological layer deformation, (middle) abrupt displacement from faulting, and (bottom) texture deformation \cite{Natali2012}.}{fig:coral-island-Natali2012}{\cite{Natali2012}'s geological storytelling method}

% These high-level approaches demonstrate the potential of sketch-based modelling not only as a geometric tool, but as a semantic interface between human intention and terrain synthesis. By abstracting terrain construction into symbolic or feature-based representations, they allow users to create rich, expressive landscapes without directly engaging with low-level geometry, making them particularly valuable for tasks involving conceptual design, education, and inverse procedural modelling.

% \midConclusion

% The methods presented in this section illustrate the diversity of sketch-based approaches for constructive terrain modelling, from curve-driven shape control to constraint-based editing and semantic abstractions. While these methods offer valuable tools for intuitive user interaction and procedural shaping, they often lack ecological grounding, multi-view integration, or the ability to produce structured data suitable for training generative models. In our work, we reinterpret and adapt elements from these approaches (dual-view sketching, curve-based region definition, and deformation fields) to support the generation of coral reef islands through a hybrid procedural and learning-based pipeline. This constructive sketch-based foundation enables us to balance user control with scalable terrain generation in data-sparse domains.

% % More recently, \cite{Liu2025} proposed a sketch-to-terrain framework that learns from a single DEM. They construct a ridge-valley graph processed by a GNN (Graph Neural Network) to preserve topological relations, and complement it with SinGAN for local detail. Their model can generate terrains of the same style as the exemplar from simple line sketches, but the geological semantics and processes (e.g., subsidence, coral keep-up) are not represented.

% % \comment{Also need to include \cite{Ketabchi2016}}


% \subsection{Deep learning}
% \label{sec:coral-island-sota-deep-learning}

% \begin{figure}
%     \includegraphics{schemaGAN_cGAN.jpg}
%     \caption[GAN and cGAN architectures]{The general structure of GAN and cGAN networks is similar: (a) a generator network $G$ is trained to take some noise $z$ as input to try to create a "realistic" output $X_{\text{fake}}$, and a discriminator network $D$ is trained in parallel to distinguish generated data from real data. (b) cGAN networks introduce a class input $c$, used by both the generator and discriminator in their inference process. In the end, only the generator is used to create new data.}
%     \label{fig:coral-island-GAN-scheme}
% \end{figure}

% Over the past decade, deep learning has revolutionised almost all areas of computer graphics and procedural content creation by learning complex, data-driven priors directly from examples. Unlike purely procedural or sketch-based methods, which rely on hand-tuned noise functions or geometric constraints, neural networks are trained to capture subtle patterns and high-frequency details without explicit programming of each effect. In terrain synthesis, this enables models to infer realistic elevation structures, textures, and region transitions from training data, even when that data is sparse or synthetic. In the context of coral reef islands, where high-resolution digital elevation models are rare, deep learning offers a way to abstract away low-level procedural rules and directly learn the mapping from semantic layouts (label maps) to plausible height fields. 
% In the following sections, we first review general Generative Adversarial Networks (GANs) and then focus on their conditional variant, cGANs, which forms the backbone of our sketch-to-terrain translation pipeline.

% \subsubsection{Generative Adversarial Networks}
% \label{sec:coral-island-sota-GAN}

% GANs, introduced in \cite{Goodfellow2014}, are a class of generative models in which two neural networks are trained in opposition (\cref{fig:coral-island-GAN-scheme}): 
% \begin{Itemize}
%     \Item{} a generator $G$ learns to produce synthetic "fake" data samples that resemble those from a target distribution given a random noise $z$, 
%     \Item{} a discriminator $D$ learns to distinguish "real" samples from those generated. 
% \end{Itemize}

% Through this adversarial process, the generator gradually improves its ability to mimic the underlying data distribution, enabling the creation of realistic outputs from random input.

% GANs have been widely adopted for image synthesis, texture generation, and data augmentation, among other tasks. In terrain modelling, they offer the potential to generate realistic landforms by learning directly from real-world data, without requiring hand-crafted procedural rules. The following works demonstrate how different GAN variants have been applied to terrain synthesis, each with its own assumptions, design trade-offs, and limitations.

% Early terrain GANs used unconditional generation, synthetising elevation maps from pure latent noise, without any spatial or semantic guidance. \citeName{WulffJensen2018} train a Deep Convolutional GAN (DCGAN) to generate realistic digital elevation models of mountainous landscapes, showing that terrain-like structures could emerge from purely data-driven learning \cite{WulffJensen2018}. The model captures local elevation statistics and allowed for latent space interpolation, enabling smooth variations across generated terrains. However, the lack of spatial conditioning make it difficult to control or constrain features such as ridges, valleys, and coastlines, resulting in landscapes that reflected training set distributions but offered no means for intentional design.

% \citeName{Spick2019} extend this approach with a Spatial GAN that generates height and texture maps jointly \cite{Spick2019}. By conditioning the generation process on spatial coordinates, their model enforced local consistency and reduced structural artefacts. This integration simplified the content pipeline by fusing geometry and appearance into a single pass. 

% Yet despite improved quality, the generation process remained fundamentally uncontrolled: users are unable to specify terrain layout, features, or semantics. As with earlier GANs, the model learned to mimic terrain distributions but could not support authoring or guided synthesis.

% %To address the lack of control inherent in purely noise-driven GANs, 
% Later works introduce multi-stage architectures in which a second conditional GAN refines or interprets the output of a first-stage generator. \citeName{Beckham2017} propose a two-step pipeline where a DCGAN generates bare terrain heightmaps from noise, and a conditional pix2pix network adds texture based on semantic cues \cite{Beckham2017}. This separation of geometry and appearance allows for basic stylisation and terrain remixing, but control for the initial terrain remains limited due to purely noise-driven GAN. Using a conceptually similar structure, \citeName{Panagiotou2020} reverse the mapping \cite{Panagiotou2020}: an unconditional GAN first synthesises aerial RGB imagery from noise, which is then passed to a cGAN trained to predict plausible DEMs. While this image-to-DEM translation enables realistic terrain reconstruction, it lacks any semantic or structural user control and relies on large paired datasets, limiting their applicability in settings such as coral reef islands, where training data are scarce. In both cases, despite the introduction of a conditional refinement stage, the generation process remains fundamentally anchored in an unconstrained latent input, offering limited authoring capability and no intuitive tool to shape landform structures.

% These methods show a shift towards using conditional models to guide terrain generation with more structure. But because they start from random noise, they offer little real control over the layout or semantic of the terrain. A natural next step is to guide the generation process directly from user-defined semantic inputs, such as sketches or label maps, to produce terrain that reflects both the training data and the user's intent.


% \subsubsection{Conditional GANs for terrain generation}
% \label{sec:coral-island-sota-cGAN}

% While traditional GANs generate data from noise, cGANs extend this concept by incorporating additional information, often called a class map or label map, to guide generation toward user-specified outcomes \cite{Mirza2014} (the $c$ class in \cref{fig:coral-island-GAN-scheme}). This makes them particularly attractive for structured content synthesis tasks, including terrain generation, where user input often defines large-scale layout and realism must emerge from learned detail. The pix2pix framework \cite{Isola2017} is the canonical cGAN formulation for image-to-image translation. It uses a U-Net generator conditioned on an input image (a sketch or a label map, for example), and a PatchGAN discriminator that evaluates realism at the patch level, thus encouraging fine detail and local consistency (\cref{fig:coral-island-pix2pix-example} presents two image-to-image translations from the same pix2pix network architecture).

% \begin{figure}
%     \includegraphics[width = 0.8 \linewidth]{example_pix2pix_facade.png}
%     \includegraphics[width = 0.8 \linewidth]{example_pix2pix_maps.png}
%     \caption[Examples of pix2pix use cases]{Example of pix2pix image-to-image translation: (top) a trained model converts semantic label maps into realistic façade images, and (bottom) constructs highly plausible aerial images from navigation map sketches. This paradigm generalises to many tasks, including terrain generation.}
%     \label{fig:coral-island-pix2pix-example}
% \end{figure}

% In terrain generation, the use of cGANs remains surprisingly rare. The most directly relevant precedent is the work of \citeName{Guerin2017}, who train a pix2pix-style cGAN to map sketched terrain features such as valleys, ridgelines, or peaks into full-resolution digital elevation models (DEMs)  \cite{Guerin2017}. Their results demonstrate that cGANs can plausibly reconstruct complex topographic forms from sparse semantic cues, offering a promising balance between user control and learned realism.
% % Similarly, \cite{Sisodia2022} apply a cGAN to generate stylised terrain heightmaps from 2D game sketch maps, validating the sketch-to-terrain pipeline. 
% \citeName{Lochner2023} updates the image-to-image process \cite{Lochner2023} by replacing the cGAN with a Denoising Diffusion Probabilistic Model \cite{Ho2020,Nichol2021}, improving the percieved quality of the resulting terrains through an iterative noising and denoising process, trained on a large DEM dataset. Diffusion models are however known to be multiple orders of magnitude slower than GANs or cGANs \cite{Xiao2022,Huang2023}, making them unusable for interactive sessions on a personal computer. 

% \citeName{Naik2022} inserts a Variational Autoencoder \cite{Kingma2022} before the cGAN phase to first compress the user input in a latent space \cite{Naik2022}, allowing to forward pass the cGAN with variations of the user input sketch or interpolate between different inputs in the 128 dimensions offered by the latent space. 

% An alternative approach from \citeName{Voulgaris2021} uses a GAN-based system that maps sparse "altitude dot" maps to plausible terrain imagery, acting as a minimal-interaction generative authoring tool \cite{Voulgaris2021}. While not strictly a cGAN, their system reflects a similar spirit: conditioning generation on lightweight user constraints. % Likewise, \cite{Panagiotou2020} and \cite{Beckham2017} train a cGAN to invert RGB satellite imagery into elevation data, framing terrain modelling as an appearance-to-geometry and geometry-to-appearance translation task. However, such image-to-DEM (and DEM-to-image) systems are heavily dependent on large paired datasets, limiting their applicability in settings such as coral reef islands, where training data are scarce.

% Despite these advances, no standard pipeline exists for generating detailed terrains from semantic layout maps using cGANs, particularly in biologically driven environments. The potential of this family of methods remains underexploited, especially when it comes to coupling user control with long-term geological plausibility.

% \midConclusion

% In our method, we address this gap by training a pix2pix cGAN to transform label maps (semantic region maps produced by procedural sketch-based modelling) into plausible coral reef island height fields. Each input map encodes key zones of an island (e.g., lagoon, reef crest, beach, island core) using categorical labels, serving as semantic constraints for terrain synthesis. The cGAN generator learns to condition elevation details on both the global layout and the implicit patterns encoded in the training data. By generating our own synthetic dataset with procedurally modelled coral reef islands (see \cref{fig:coral-island-difficulties-dataset}), we overcome the shortage of labelled elevation data and enforce geological coherence via data generation design. The choice of the pix2pix architecture is tailored by its short training and inference time and the availability of well-maintained implementations that support reproducibility, but any other conditional architecture could be substitued such as pix2pixHD or BicycleGAN \cite{Wang2018,Zhu2018NN}.

% This setup offers several key advantages: 
% \begin{Itemize}
%     \Item{} respecting user-defined structure while allowing the generator to introduce realistic variation,
%     \Item{} removing procedural biases (radial symmetry and fixed island typologies),
%     \Item{} enabling the generation of irregular, non-circular landforms while implicitly modelling geological processes such as subsidence and coral accretion through training-time priors.
% \end{Itemize}


% \begin{figure}
%     \autofitgraphics[]{examples-going-more-complex.pdf}
%     \caption[Examples of procedural region maps used for training a cGAN]{Examples of procedural region maps used for training: left to right, canonical island, off-centre island, elongated shapes, and multi-island scenes. These maps serve as semantic inputs to our cGAN.}
%     \label{fig:coral-island-difficulties-dataset}
% \end{figure}


% \bigConclusion

% In this context, conditional GANs emerge as a powerful yet underutilised tool for terrain modelling. By training on procedurally generated coral reef island data, we show that sketch-conditioned learning can effectively bridge the gap between high-level user intent and geologically plausible terrain synthesis. We note that few-shot or one-shot learning methods are emerging such as \citeName{Liu2025}'s Graph Neural Network usage, possibly solving the dataset scarcity problem, at the cost of methods falling out of the interactive-time scope with multiple seconds at inference time \cite{Liu2025}.

